%%% Abstract file — content only, no \chapter or formatting commands.
%%% UF-F15: must be 350 words or fewer (source: bundled template's own
%%% abstractFile.tex states this cap).
%%% The class file's \AtBeginDocument block builds the abstract page
%%% header (degree info, chair, major) automatically — the student just
%%% provides the abstract prose below.

Graduate schools enforce strict formatting rules on doctoral
dissertations, but the rules are scattered across web pages, PDF
guides, and template documentation. Students who already compile their
dissertation locally still face a tedious manual checklist before
submission, often performed under time pressure near the term deadline.
A formatting mistake caught after the first submission deadline can
delay graduation by a full semester.

This dissertation presents a lightweight static analysis tool that
validates LaTeX-based dissertations against the University of Florida
Graduate Editorial Office's formatting specification. The tool operates
in two complementary layers. A source layer inspects the LaTeX project
for overrides of the official UF template's defaults — for example, a
manually inserted \texttt{\textbackslash pagenumbering\{roman\}} that
would violate the new template's arabic-throughout rule. A PDF layer
inspects the rendered output for properties the source cannot
predetermine: actual margins per page, embedded font families, page
numbering scheme as rendered, and tagged-structure presence for
accessibility compliance. The two layers are designed to complement
each other: source-layer findings carry precise location data; PDF-layer
findings carry rendered-truth confidence.

The analyzer accepts project archives, project directories, or compiled
PDFs as input. Output is a grouped, severity-tiered report with each
finding citing the originating UF rule. A demo dissertation that
satisfies every rule serves as both regression fixture and pedagogical
reference for new graduate students learning the template.

This work focuses on doctoral dissertations using the Fall 2025+ UF
template. Master's theses and the older template (accepted through
Summer 2026) are out of scope for the initial release. Evaluation on a
corpus of public UF dissertations shows the source-layer checks alone
recover the majority of formatting violations recorded in the Editorial
Office's published rejection patterns, with the PDF layer adding
defense-in-depth against template bypass scenarios.
